% ATTEMPT_STATUS: unresolved
% TOP_PROBLEM: {"problemId": "problem.yang-mills-existence-and-mass-gap", "canonicalTitle": "Yang-Mills Existence and Mass Gap", "releaseRank": 3, "primaryDomain": "mathematical_physics", "primaryDomainLabel": "Mathematical physics", "displayStatus": "open", "releaseStatus": "open"}
% !TeX program = lualatex
\documentclass[11pt]{article}
\usepackage[margin=25mm]{geometry}
\usepackage{fontspec}
\setmainfont{Latin Modern Roman}
\usepackage{amsmath,amssymb,mathtools}
\usepackage{unicode-math}
\setmathfont{Latin Modern Math}
\usepackage{xurl,hyperref}
\hypersetup{colorlinks=true,urlcolor=blue}
\setcounter{secnumdepth}{0}
\setlength{\parindent}{0pt}
\setlength{\parskip}{5pt}
\setlength{\emergencystretch}{3em}
\begin{document}
\section*{\texorpdfstring{Yang-Mills Existence and Mass Gap}{Yang-Mills Existence and Mass Gap}}\label{3}
\subsection{Definitions and mathematical statement}
A classical $G$-connection has curvature $F_A=dA+A\wedge A$ and Yang--Mills action proportional to $\int_{{\mathbb{R}}^4}\langle F_A,F_A\rangle\,dx$. The question concerns a rigorous, nontrivial \emph{quantum} theory associated with this action, not only classical solutions or a formal functional integral. In its vacuum representation, let $\Omega$ be the vacuum and $H\ge0$ the self-adjoint generator of time translations, with $H\Omega=0$. A positive finite mass gap means
\[
 0<m:=\inf\bigl(\operatorname{spec}(H)\setminus\{0\}\bigr)<\infty.
\]
For every compact simple $G$, construct such a theory with the local observable fields, positivity, covariance, locality, and short-distance conditions required in Jaffe--Witten, Sections 3--4. Those axioms are part of the target and are not replaced by the displayed spectral condition alone.
\par\smallskip\textit{Formulation and concept references:} \href{https://www.claymath.org/wp-content/uploads/2022/06/yangmills.pdf}{[S1]}.

\subsection{Short English statement}
Can every compact simple gauge symmetry produce a rigorously defined four-dimensional Yang--Mills quantum theory with a nonzero energy gap above its vacuum?
\subsection{Research attempt}

\paragraph{Scope and source audit (27 September 2026).}
This attempt treats every compact simple gauge group separately; constants may
 depend on the group. It requires a nontrivial four-dimensional continuum
 theory with the observable fields and ultraviolet properties in [S1]. A
 finite lattice, a classical Yang--Mills solution, or a spectral estimate for
 one gauge-fixed propagator does not meet those requirements.

The catalogue's recent references need explicit qualification. Nair's [S3]
 describes an outline for \emph{three}-dimensional Yang--Mills theory. The arXiv
 record for Jacobsen's [S4] marks version 2 as withdrawn by arXiv Admin on
 13 June 2025, with a research-content quality notice. Neither source is used
 here as a proved four-dimensional construction. The lattice positivity
 framework of Osterwalder--Seiler [E1] supplies useful finite-cutoff structure;
 Chatterjee's own lecture notes [E2] distinguish fixed-coupling lattice
 estimates from the weak-coupling continuum problem. The deductions below
 are elementary spectral and limiting arguments, not claims of new results
 in constructive field theory.

\paragraph{Attack map and the first object that actually exists.}
The constructive route has three tasks: construct renormalized Euclidean
 correlations, obtain estimates uniform as the cutoff and volume are removed,
 and reconstruct the physical Hilbert space. A direct Hamiltonian route would
 instead require a well-defined renormalized operator and its spectral gap.
 I pursue the first route, while testing whether finite-cutoff positivity or
 a single massive correlation channel could close its missing step.

Choose a faithful unitary representation $\rho:G\to U(d_\rho)$ and a finite
 hypercubic lattice. Assign $U_e\in G$ to one orientation of each edge, put
 $U_{\bar e}=U_e^{-1}$, and let $U_p$ be the ordered plaquette product. For
 $\beta\geq0$ define
\[
 S_\beta(U)=\beta\sum_p
 \left(1-d_\rho^{-1}\operatorname{Re}\operatorname{Tr}\rho(U_p)\right),
 \qquad d\mu_\beta=Z_\beta^{-1}e^{-S_\beta(U)}\prod_e dU_e.
\]
Here $dU_e$ is normalized Haar measure. Since every summand is in $[0,2]$,
\[
 e^{-2\beta N_p}\leq Z_\beta\leq1.
\]
Thus this measure exists without perturbation theory. Haar invariance and
 invariance of the trace under conjugation prove gauge invariance. Bounded
 observables have finite moments. These facts are exact, but the bound on
 $Z_\beta$ degenerates with the plaquette count and gives no compactness for
 renormalized local fields as the lattice spacing decreases.

Reflection positivity for the appropriate Wilson lattice construction is a
 separate theorem [E1]. It is not a consequence of the ordinary positivity
 of $d\mu_\beta$: reflection positivity requires
 $\langle\overline{F(\theta U)}F(U)\rangle\geq0$ for positive-time observables
 $F$. I use the resulting transfer-operator picture only where its
 hypotheses are explicitly assumed below. No global gauge fixing is needed
 for the finite measure or the spectral argument.

\paragraph{Lemma 3.1: what exponential decay really proves.}
Let $H\geq0$ be self-adjoint on a Hilbert space, let $\Omega$ be a unit vector
 with $H\Omega=0$, and let $\mathcal D$ be a dense linear subspace of
 $\Omega^\perp$. Suppose there is one $m_0>0$ such that for every
 $v\in\mathcal D$ there are finite $K_v$ and $t_v$ with
\begin{equation}\label{ym:decay}
 0\leq\langle v,e^{-tH}v\rangle\leq K_v e^{-m_0t}
 \quad(t\geq t_v).
\end{equation}
Then $\ker H=\mathbb C\Omega$ and
 $\operatorname{spec}(H)\subset\{0\}\cup[m_0,\infty)$.

\emph{Proof.}
Write $E_H$ for the spectral projection measure and
 $\mu_v(B)=\langle v,E_H(B)v\rangle$. For $0\leq b<m_0$,
\[
 e^{-bt}\mu_v([0,b])
 \leq\int_{[0,\infty)}e^{-t\lambda}\,d\mu_v(\lambda)
 \leq K_v e^{-m_0t}.
\]
Letting $t\to\infty$ gives $\mu_v([0,b])=0$. Equivalently,
 $E_H([0,b])v=0$. Boundedness of this projection extends the equality to
 $\Omega^\perp$. Taking a countable increasing sequence $b\uparrow m_0$
 excludes all spectrum in $[0,m_0)$ on that subspace. Since $H\Omega=0$
 and $H$ is self-adjoint, $\Omega^\perp$ reduces $H$. The conclusions follow.
\hfill$\square$

The common exponent matters; the constants need not be uniform in $v$.
 Diagonal matrix elements are nonnegative by the spectral theorem, so this
 argument does not silently discard cancellations in signed correlations.
 Conversely, the asserted gap and unique vacuum imply
 \eqref{ym:decay} with $K_v=\|v\|^2$ for every $v\perp\Omega$.

There is also a nontriviality check. If a nonzero $v\perp\Omega$ survives,
 then $C_v(0)=\|v\|^2>0$ and $C_v(t)=\langle v,e^{-tH}v\rangle>0$ for every
 finite $t>0$. The latter is strict because $e^{-t\lambda}>0$ at every finite
 spectral value. The spectral gap $m$ consequently satisfies
\begin{equation}\label{ym:upper}
 m_0\leq m\leq-\frac1t\log\frac{C_v(t)}{C_v(0)}<\infty.
\end{equation}
Indeed $C_v(t)\leq e^{-mt}C_v(0)$. This upper bound concerns the bottom of
 the excitation spectrum, not the existence of an isolated particle state.

\paragraph{A precise conditional continuum reduction.}
Take lattice spacings $a_j\downarrow0$ and spatial side lengths $L_j$ with
 $a_jL_j\to\infty$. Time must also become unbounded in physical units;
 a fixed periodic time circle would describe a thermal problem. Suppose the
 following properties can be proved for renormalized gauge-invariant fields.

\begin{enumerate}
\item Their complete family of smeared Euclidean correlations converges and
 satisfies the hypotheses needed for Osterwalder--Schrader reconstruction,
 including regularity, reflection positivity, Euclidean covariance, and
 the required compatibility conditions. The reconstructed observable fields
 have the short-distance behavior and field content required by [S1].
\item Positive-time observable polynomials, after subtracting their vacuum
 components and quotienting null vectors, give a dense family
 $v_F\in\Omega^\perp$. Their lattice representatives $v_{F,j}$ have diagonal
 transfer correlations $C_{F,j}(n)$ such that for every fixed $t>0$,
\[
 C_{F,j}(\lfloor t/a_j\rfloor)
 \longrightarrow\langle v_F,e^{-tH}v_F\rangle.
\]
The limiting norm is finite, and at least one such $v_F$ is nonzero.
\item There exists $m_0(G)>0$ such that, for each $F$, constants $K_F,t_F$
 independent of $j$ satisfy
\begin{equation}\label{ym:uniform}
 0\leq C_{F,j}(n)\leq K_F e^{-m_0(G)a_jn}
 \quad\text{whenever }a_jn\geq t_F.
\end{equation}
\end{enumerate}

For any fixed $t>t_F$, passage to the limit in \eqref{ym:uniform} proves
 \eqref{ym:decay}. Lemma 3.1 then gives the mass gap and uniqueness of the
 vacuum, while the surviving vector and \eqref{ym:upper} give a finite mass.
 This proves the stated implication. It does \emph{not} prove any of the
 three assumptions for four-dimensional Yang--Mills theory. In particular,
 reconstructing an arbitrary massive field theory would not identify it as
 the Yang--Mills theory specified by the catalogue.

These assumptions separate two logically different limiting operations.
 Compactness of bounded Wilson-loop expectations could produce subsequential
 numbers, but local fields require normalization, smearing, and control of
 coincident-point singularities. Also, estimates established only after taking
 volume to infinity at each fixed $a$ need cutoff-uniform constants before
 they imply the displayed diagonal limit. Neither order of limits is licensed
 by boundedness of a finite partition function.

\paragraph{Failed shortcut A: a lattice gap without physical units.}
If a positive normalized transfer operator has largest eigenvalue $1$ and
 next spectral edge $\lambda_1(a)<1$, its dimensionless gap is
 $\delta(a)=-\log\lambda_1(a)$. A step advances physical time by $a$, so
\[
 H_a=-a^{-1}\log T_a,\qquad m_a=\frac{\delta(a)}a.
\]
On the orthogonal complement of the vacuum, the three diagonal examples
\[
 \lambda_1(a)=e^{-a^2},\qquad e^{-ma},\qquad e^{-c}
 \quad(m,c>0)
\]
all have strictly positive gaps at every cutoff, but their physical gaps
 tend respectively to $0$, $m$, and $\infty$. Consequently finite-cutoff
 positivity cannot determine the continuum answer. The last example has
 $C_a(t)=e^{-c\lfloor t/a\rfloor}$ for a unit excited vector. Its limit is
 zero at every $t>0$, yet is one at $t=0$; it cannot be the correlation of a
 fixed unit vector under a strongly continuous continuum semigroup.

This explains why a fixed strong-coupling estimate, even if excellent in
 lattice units, is insufficient. The coupling trajectory, renormalized
 observables, and a finite physical scale must all be controlled together.
 The distinction agrees with the weak-coupling continuum discussion in [E2].
 It does not assert that every strong-coupling method must fail; it identifies
 the extra estimate that would be needed to use one here.

\paragraph{Failed shortcut B: one observable misses light sectors.}
Consider
\[
 \mathcal H=\mathbb C\Omega\oplus\mathbb Cv
       \oplus\ell^2(\mathbb N),\qquad
 H\Omega=0,\quad Hv=v,\quad He_n=n^{-1}e_n.
\]
This is a bounded positive self-adjoint operator with a unique zero
 eigenvector and no positive gap. Nevertheless
 $\langle v,e^{-tH}v\rangle=e^{-t}$. Thus an exactly massive channel, not
 just a numerical approximation to one, cannot prove the target. Each
 $e_n$ also has exponential decay, but with exponent $1/n$; allowing the
 exponent to depend on the observable would invalidate Lemma 3.1. Physical
 local observables must generate a dense family, and the decay exponent must
 be common to that family.

A further limit stress test takes
 $\mathcal H_j=\mathbb C\Omega\oplus\operatorname{span}(e_1,\ldots,e_j)$
 with energies $1,1/2,\ldots,1/j$. Every approximant has a nontrivial vacuum
 gap, while the embedded strong limit is gapless. A volume-dependent gap
 tending to zero is therefore not repaired merely by convergence of the
 Hamiltonians on finite-support vectors.

\paragraph{Verification and exact remaining gap.}
The proofs above were checked against a zero-dimensional excitation space,
 extra zero modes, a continuum of positive spectral values, and constants
 depending on the observable. The nonzero-vector requirement excludes the
 vacuum-only theory. The projection argument excludes extra zero modes and
 does not require eigenvectors above the vacuum. The units calculation uses
 physical time $an$, rather than replacing it by the number of steps.
 The diagonal examples give explicit counterexamples to each weakened
 spectral assertion. They are Hilbert-space models, not counterexamples to
 Yang--Mills existence.

The first unproved analytic task on this route is to produce one
 renormalized continuum-limit family with the full required field content
 and a nonzero physical observable norm. Even conditional on that
 construction, \eqref{ym:uniform} for a dense observable family remains
 unproved. The gap lemma solves only the spectral inference after those
 inputs are supplied.

Three focused next steps are to establish uniform bounds for a chosen
 family of smeared curvature-polynomial correlations; to track the
 transfer-correlation exponent in physical units along the same cutoff
 trajectory; and to exhibit a normalized observable with a nonzero limiting
 connected correlation. Each must retain control as the physical volume
 grows. A hypothetical obstruction could arise from low-energy states missed
 by the tested observables, or from loss of all nonvacuum vectors in the
 continuum limit. The examples show how to detect these failures without
 mistaking them for a disproof of the original conjecture.

\paragraph{Final status and provenance.}
\textbf{UNRESOLVED.} This attempt proves a conditional spectral criterion
 and a finite-mass test, and supplies explicit counterexamples to several
 tempting shortcuts. It constructs no four-dimensional continuum
 Yang--Mills theory. Written by GPT-6 Astra Ultra from the numbered target,
 using direct primary-source checks and an adversarial audit of cutoff,
 density, and nontriviality assumptions. No theorem-prover verification or
 numerical lattice simulation is claimed.

\subsection{Sources}
\begin{itemize}
\item[S1]Arthur Jaffe and Edward Witten, Quantum Yang-Mills Theory, official Clay problem description.
\url{https://www.claymath.org/wp-content/uploads/2022/06/yangmills.pdf}.
\textit{Formulation source. Consulted 24 September 2026: Sections 3--4, especially page 6.}
\item[S2]Yang-Mills \& the Mass Gap.
\url{https://www.claymath.org/millennium/yang-mills-the-maths-gap/}.
\textit{Further reference; not a proof endorsement.}
\item[S3]Towards a Proof of Mass Gap in 3d Yang-Mills Theory — V. P. Nair.
\url{https://arxiv.org/abs/2608.10133}.
\textit{Further reference; not a proof endorsement.}
\item[S4]A Constructive Proof of Existence and Mass Gap for Pure SU(3) Yang-Mills in Four-Dimensional Space-Time — D. C. Jacobsen.
\url{https://arxiv.org/abs/2506.00284}.
\textit{Further reference; not a proof endorsement.}
\item[E1]Konrad Osterwalder and Erhard Seiler, \emph{Gauge field theories on a lattice}, Annals of Physics \textbf{110} (1978), 440--471.
\url{https://doi.org/10.1016/0003-4916(78)90039-8}.
\textit{Primary lattice-positivity reference; publisher abstract checked 27 September 2026. The complete constructive argument was not independently reverified.}
\item[E2]Sourav Chatterjee, \emph{Yang--Mills on the lattice: New results and open problems}, author-hosted lecture slides, especially slides 14--19 and 33--35.
\url{https://souravchatterjee.su.domains/beam-owps-trans.pdf}.
\textit{Author's statements about lattice estimates and their continuum scope; consulted 27 September 2026.}
\end{itemize}
\end{document}
